This chapter follows Section~4 of the paper, \emph{Particle dynamics and absence of
exponential stretching} (\texttt{sec:particle}).  It has two halves: the
particle flow of the global velocity, written down in closed form, with its linear Lipschitz
bound and its zero topological entropy (Proposition~\ref{prop:particle-flow}); and the ideal
induction group built from the Cauchy formula, with its two-sided polynomial bounds and its
vanishing exponential growth rate (Corollary~\ref{cor:ideal-growth}).  Together they give
the two dynamical clauses of Theorem~\ref{thm:main}, and rather more
(Proposition~\ref{prop:main-lagrangian}).

The predicates are stated in
\texttt{FastDynamo/Dynamics/ParticleFlow.lean}; the flow is constructed in closed form, cell
by cell, in \texttt{FastDynamo/Dynamics/Flow/*}; the ideal group is constructed from the
Cauchy formula in \texttt{FastDynamo/Dynamics/Ideal/*}; and
\texttt{FastDynamo/Dynamics/MainLagrangianClauses.lean} assembles them at the constructed
velocity (\texttt{mainTheoremLagrangian\_holds}), which discharges the hypothesis of
\texttt{mainTheoremComplete\_of\_lagrangian} and of \texttt{ideal\_resistive\_gap}.  Almost
no ODE theory is used: Picard--Lindel\"of, integral curves and Rademacher are avoided
altogether by writing the flow down.  The steps of the paper's \emph{proof} of Corollary~4.2
(\texttt{cor:ideal-growth}) are formalised: the variational equation
\texttt{eq:flow-variational-equation} (Lemma~\ref{lem:dynamics:variational-equation}), the
weak form of the ideal equation and the fact that the Cauchy formula solves it
(Definition~\ref{def:dynamics:ideal-weak-form}, Lemma~\ref{lem:dynamics:ideal-weak-existence}),
and the uniqueness of weak solutions in $\C([-T,T];\Lp^2)$, which makes $\mathcal T_0$
\emph{the} solution group (Lemma~\ref{lem:dynamics:ideal-weak-uniqueness}).  The last section
is Remark~1.7 of the paper, the convergence of the resistive evolution to the ideal one as
$\eps\downarrow0$ (Theorem~\ref{thm:vanishing-diffusivity}).

\section{The particle flow and its entropy}

\begin{definition}[Particle flow and Lipschitz constants]
  \label{def:dynamics:particle-flow}
  \lean{FastDynamo.literalTorusPoint, FastDynamo.IsLiteralTorusLift,
    FastDynamo.IsParticleFlow,
    FastDynamo.ScaleFamily.TorusCellPlacement.literalFlow,
    FastDynamo.ScaleFamily.TorusCellPlacement.literalFlowLift,
    FastDynamo.mainLiteralFlow,
    FastDynamo.IsParticleFlow.unique,
    FastDynamo.ScaleFamily.TorusCellPlacement.literalFlow_unique,
    FastDynamo.mainLiteralFlow_unique,
    FastDynamo.existsUnique_isParticleFlow_mainLiteralVelocity}
  \leanok
  \uses{def:velocity:global-velocity, lem:global-velocity}
  The \emph{particle flow} of the global velocity $u $ is the family
  $\Phi_t:\T^3\to\T^3 $ solving
  \[
    \frac{\dd}{\dd t}\Phi_t(x)=u(\Phi_t(x)),
    \qquad
    \Phi_0(x)=x.
  \]
  Existence and uniqueness for every $t\in\R $ follow from the global Lipschitz regularity
  of $u $ (Lemma~\ref{lem:global-velocity}), and $\Phi_t^{-1}=\Phi_{-t} $.  For a
  Lipschitz map $F $, $\operatorname{Lip}(F) $ denotes its optimal Lipschitz constant.
\end{definition}

There is both a Lean statement and a Lean flow.  \texttt{IsParticleFlow v} $\Phi$ says
that for every $x$ there is a single curve $\gamma\colon\R\to\R^3$ with
$\Phi_t(x)=\pi(\gamma(t))$ and $\gamma'(t)=v(\Phi_t(x))$, together with $\Phi_0=\id$; the
inverse law $\Phi_{-t}\circ\Phi_t=\id$ (used in the proof of Proposition~4.1 of the paper
as $\Phi_t^{-1}=\Phi_{-t}$) is a separate conjunct of
\texttt{MainLagrangianClauses}, and the full group law $\Phi_s\circ\Phi_t=\Phi_{s+t}$ is proved
for the periodised data from which the flow is built
(\texttt{PeriodizedPlaneFlowData.literalFlow\_comp}) but is not part of any clause.  The paper's own
device of working with periodic lifts is
\texttt{IsLiteralTorusLift}.  The flow itself is
\texttt{ScaleFamily.TorusCellPlacement.literalFlow}
(\texttt{FastDynamo/Dynamics/Flow/TorusFlow.lean}), with periodic lift
\texttt{literalFlowLift}, and \texttt{mainLiteralFlow} is its instance at the constructed
velocity.  Its \emph{existence} is built \emph{in closed form}, not by an existence theorem:
the library never invokes Picard--Lindel\"of or integral curves, where the paper obtains
existence and uniqueness of the flow from the global Lipschitz bound of
Lemma~3.1 of the paper (opening paragraph of Section~4) and then derives the closed formula
\texttt{eq:explicit-particle-flow} for that flow.  The closed formula itself is the
paper's -- the transverse rotation of
\texttt{Global/PairChartRotation.lean}, the radial angular and axial profiles of
\texttt{Dynamics/Flow/CellProfiles.lean}, the single-cell map of
\texttt{Dynamics/Flow/CellFlowMap.lean}, the sum over the accumulating cells and its
$\Z^2$-periodisation in \texttt{Dynamics/Flow/PlaneFlow\{,Regularity\}.lean}, and the
descent through the chart in \texttt{Dynamics/Flow/TorusFlow.lean}.  That the closed
formula solves the ODE is Proposition~\ref{prop:particle-flow}.

\emph{Uniqueness} is proved, by Gronwall's inequality.  The paper's
``the unique particle flow'' (opening paragraph of Section~4) is
\texttt{IsParticleFlow.unique}
(\texttt{FastDynamo/Dynamics/Flow/ParticleFlowUniqueness.lean}): for a Lipschitz $v$ there is at
most one $\Phi$ with $\texttt{IsParticleFlow}\ v\ \Phi$.  The argument runs on the universal
cover, in the spirit of the paper's periodic lifts.  Two flows
give, through a fixed point $x$, two lifted trajectories $\gamma,\gamma'$ solving the same
\emph{autonomous} ODE $\dot r=\tilde v(r)$ with $\tilde v=v\circ\pi$, which is Lipschitz with
the same constant because $\pi$ is $1$-Lipschitz
(\texttt{lipschitzWith\_literalTorusPoint}).  Their starting points agree modulo $\Z^3$, so the
translate $\gamma'+(\gamma(0)-\gamma'(0))$ solves the same equation with the same initial value,
and Mathlib's \texttt{dist\_le\_of\_trajectories\_ODE} (Gronwall with $\delta=0$), applied
forwards and to the time-reversed system, identifies the two curves; projecting back gives
$\Phi=\Phi'$.  Hence \texttt{literalFlow} is \emph{the} flow at every admissible parameter
tuple (\texttt{literalFlow\_unique}), and \texttt{mainLiteralFlow} is the paper's
$\Phi_t$ at the constructed velocity (\texttt{mainLiteralFlow\_unique},
\texttt{existsUnique\_isParticleFlow\_mainLiteralVelocity}).  Nothing downstream needs it:
$\Phi_{-t}\circ\Phi_t=\id$ is a clause of \texttt{MainLagrangianClauses} in its own right, not a
consequence of uniqueness.

One deviation, forced by the geometry: the chart-side plane flow is \emph{not}
$\Z^2$-equivariant (its displacement is supported in the placement chart, which is
disjoint from its nontrivial integer translates), exactly as the velocity is not, so the
object that descends to the torus is the periodisation
\texttt{planePeriodizedTransverseFlow} built with the same
\texttt{UnitBoxPeriodization.periodization} the velocity already uses.  On the placement
chart the two agree, so no constant is lost.

\begin{definition}[Topological entropy]
  \label{def:dynamics:topological-entropy}
  \lean{FastDynamo.HasZeroTopologicalEntropy,
    FastDynamo.coverEntropy_eq_zero_of_lipschitz_iterate}
  \leanok
  Let $F $ be a continuous self-map of a compact metric space $(X,d) $, and write
  $h_{\mathrm{top}}(F) $ for its topological entropy.  The paper fixes no display for
  it: it refers to [Hasselblatt--Katok, \emph{Principal structures}, Section~2.5e]
  (Section~1.2 of the paper, before Theorem~1.2).  That section introduces $h_{\mathrm{top}} $ through the dynamical metrics
  $d_m(x,y)=\max_{0\leq j<m}d(F^j(x),F^j(y)) $ and the exponential growth rate in $m $ of
  the minimal cardinality of an $(m,\rho) $-cover, equivalently of the maximal cardinality
  of an $(m,\rho) $-separated set; on a compact metric space the cover, spanning and
  separated forms are one and the same invariant.
\end{definition}

\textbf{Which invariant.}  The Lean predicate \texttt{HasZeroTopologicalEntropy F} is
$\texttt{Dynamics.coverEntropy}\ F\ \texttt{univ}=0$, Mathlib's \emph{cover} entropy.  On
the compact metric space $\T^3$ the proof reduces it through
\texttt{Metric.uniformity\_basis\_dist} (inside
\texttt{coverEntropy\_eq\_zero\_of\_lipschitz\_iterate}) to the growth rate of the minimal
$(m,\rho)$-cover, so the quantity proved to vanish is the Bowen--Dinaburg invariant of the
cited section, not a weaker substitute.  Mathlib's own equivalence with the separated (net)
form, \texttt{Dynamics.coverEntropy\_eq\_iSup\_netEntropyEntourage}, exists in the
pinned Mathlib but is not imported here, because no clause needs it
(\texttt{Mathlib/Dynamics/TopologicalEntropy/NetEntropy.lean} is not built in this tree).

\textbf{The proof route differs.}  The paper (last paragraph of the proof of
Proposition~4.1) combines two facts quoted
from Hasselblatt--Katok: the power rule $h_{\mathrm{top}}(F^m)=m\,h_{\mathrm{top}}(F)$
(Section~2.5f(3)) and the Lipschitz bound
$h_{\mathrm{top}}(F)\leq3\log\operatorname{Lip}(F)$ (Proposition~2.5.5), applied to
$F=\Phi_s$, $F^m=\Phi_{ms}$ and $\operatorname{Lip}(\Phi_{ms})\leq C_{\rm L}(1+m\abs s)$.
Mathlib has neither fact -- there is no entropy power rule in
\texttt{Mathlib/Dynamics/TopologicalEntropy/} and no entropy bound for a Lipschitz self-map
-- so the development proves the \emph{combination} directly instead of its two factors.
\texttt{coverEntropy\_eq\_zero\_of\_lipschitz\_iterate}
(\texttt{FastDynamo/Analysis/EntropyCriterion.lean}) says that if every iterate $F^{[m]}$ is
$C(1+m)$-Lipschitz on $\T^3$ then $\texttt{coverEntropy}\,F\,\texttt{univ}=0$, and
\texttt{lipschitzWith\_literalFlow\_iterate} supplies that hypothesis from the group law
$\Phi_s^{[m]}=\Phi_{ms}$ and the Lipschitz bound of
Proposition~\ref{prop:particle-flow}.  Its proof is the classical Bowen--Dinaburg covering
argument that underlies the Lipschitz bound of Hasselblatt--Katok:
the $N^3$-point grid of \texttt{Analysis/TorusNet.lean} is a $1/(2N)$-net of $\T^3$, so it
is a dynamical $(m,\rho)$-cover for $N=\lceil C(1+m)/\rho\rceil+1$, the covering number
grows polynomially in $m$, and the exponential growth rate of a polynomial is zero.  The
input is the paper's (the linear Lipschitz bound and the group law); only the two quoted
theorems are bypassed, because formalising either of them would be strictly more work than
the combination.

\begin{proposition}[Explicitly solvable particle dynamics and zero entropy]
  \label{prop:particle-flow}
  \lean{FastDynamo.IsParticleFlow, FastDynamo.HasZeroTopologicalEntropy,
    FastDynamo.MainLagrangianClauses,
    FastDynamo.ScaleFamily.TorusCellPlacement.isParticleFlow_literalFlow,
    FastDynamo.ScaleFamily.TorusCellPlacement.forall_literalFlow_neg_literalFlow,
    FastDynamo.ScaleFamily.TorusCellPlacement.lipschitzWith_literalFlow,
    FastDynamo.ScaleFamily.TorusCellPlacement.exists_lipschitzWith_literalFlow,
    FastDynamo.ScaleFamily.TorusCellPlacement.coverEntropy_literalFlow_eq_zero,
    FastDynamo.ScaleFamily.TorusCellPlacement.forall_hasZeroTopologicalEntropy_literalFlow,
    FastDynamo.lipConst,
    FastDynamo.ScaleFamily.TorusCellPlacement.lipConst_literalFlow_le,
    FastDynamo.ScaleFamily.TorusCellPlacement.inv_flowLipschitzBound_le_lipConst_literalFlow,
    FastDynamo.ScaleFamily.TorusCellPlacement.tendsto_abs_inv_mul_log_lipConst_literalFlow,
    FastDynamo.tendsto_abs_inv_mul_log_lipConst_mainLiteralFlow,
    FastDynamo.ScaleFamily.TorusCellPlacement.literalFlow_eq_cell_of_mem_bufferDisk,
    FastDynamo.mainLiteralFlow_eq_cell_of_mem_bufferDisk,
    FastDynamo.ScaleFamily.TorusCellPlacement.literalFlow_eq_self_of_forall_not_mem_bufferDisk,
    FastDynamo.ScaleFamily.TorusCellPlacement.literalFlow_eq_self_of_accumulationPoint,
    FastDynamo.mainLiteralFlow_eq_self_of_accumulationPoint,
    FastDynamo.ScaleFamily.TorusCellPlacement.cellAngularSpeed_cellCenter,
    FastDynamo.ScaleFamily.TorusCellPlacement.cellAxialSpeed_cellCenter}
  \leanok
  \uses{def:dynamics:particle-flow, def:dynamics:topological-entropy, def:velocity:global-velocity, def:velocity:scale-family, def:velocity:placement}
  Let $u $ be the global velocity field.  There is a constant $C_{\rm L}\geq1 $ with
  \[
    \operatorname{Lip}(\Phi_t)+\operatorname{Lip}(\Phi_t^{-1})
    \leq C_{\rm L}(1+\abs t)
    \qquad(t\in\R).
  \]
  In particular $\lim_{\abs t\to\infty}\abs t^{-1}\log\operatorname{Lip}(\Phi_t)=0 $, and
  every time map has zero topological entropy:
  \[
    h_{\mathrm{top}}(\Phi_t)=0
    \qquad(t\in\R).
  \]

  The flow is moreover explicit.  With $\mathsf R_\alpha $ the rotation by the angle
  $\alpha $ in the horizontal plane and, for the $n $-th cell,
  \[
    \omega_n(r)=\Omega\chi_h\!\left(\frac r{\ell_n}\right),
    \qquad
    g_n(r)=\ell_nU\chi_h\!\left(\frac r{\ell_n}\right)
    +a_n\vartheta\!\left(\frac r{R_n}\right),
  \]
  one has, for $y\in D(p_n,R_n) $, $r=\abs{y-p_n} $ and $x_3\in\T $,
  \begin{equation}
    \Phi_t(y,x_3)
    =\left(p_n+\mathsf R_{t\omega_n(r)}(y-p_n),\ x_3+tg_n(r)\right),
    \label{eq:dynamics:explicit-flow}
  \end{equation}
  the second component modulo $1 $.  At $y=p_n $ the horizontal component is $p_n $ while
  the axial component translates at the speed $g_n(0)=\ell_nU+a_n $.  Every point outside
  the cells, on their boundaries, and of the accumulation circle
  $\{p_\infty\}\times\T $ is fixed.
\end{proposition}

Proved, in every assertion.  The ODE clause is
\texttt{isParticleFlow\_literalFlow}, the inverse clause
$\Phi_{-t}\circ\Phi_t=\id$ is \texttt{forall\_literalFlow\_neg\_literalFlow}, the Lipschitz
bound is \texttt{exists\_lipschitzWith\_literalFlow} with the explicit constant
$C_{\rm L}=\texttt{flowCL}=4(1+s)$, $s=\texttt{flowLipschitzSlope}$, the log-Lipschitz limit
is \texttt{tendsto\_abs\_inv\_mul\_log\_lipConst\_literalFlow}, and the entropy clause is
\texttt{coverEntropy\_literalFlow\_eq\_zero}.  All of them except the log-Lipschitz limit are
conjuncts of the Lean predicate \texttt{MainLagrangianClauses}, and
Proposition~\ref{prop:main-lagrangian} is that predicate at the constructed velocity; the
log-Lipschitz limit is a theorem beside it
(\texttt{tendsto\_abs\_inv\_mul\_log\_lipConst\_mainLiteralFlow}), because no clause of
Theorem~\ref{thm:main} consumes it.

Four statements are rendered differently from the paper.  First,
$\operatorname{Lip}$ is \texttt{FastDynamo.lipConst}, the infimum of the set of admissible
Lipschitz constants.  That this infimum is itself admissible is \emph{not} proved, and is not
needed: the upper bound is an \texttt{sInf} below an admissible constant and the lower bound
holds for every admissible one, so the squeeze goes through without it.  The bundle's own
Lipschitz clause quantifies existentially over admissible constants instead, which is the
same statement, the optimal constant being the least admissible one.  Second, $\abs t\to\infty$
is read as \texttt{Filter.cocompact}~$\R$, the standard Mathlib reading of that limit.
Third, \eqref{eq:dynamics:explicit-flow} is realised twice: as the \emph{definition} of
\texttt{literalFlow} -- the horizontal component is \texttt{planeRotate} at the angle
$t\,\omega_n(r)$ about the cell centre and the axial component the shift by $t\,g_n(r)$,
with the radius a constant of the motion
(\texttt{cellTransverseFlow\_mem\_bufferDisk}), so that the closed form is what the ODE
verification consumes -- and, as a theorem about the torus flow,
\texttt{literalFlow\_eq\_cell\_of\_mem\_bufferDisk}, whose hypotheses pin $y$ as a
fundamental-box representative of the horizontal block of $x$; that is what makes the
chart-side cell formula visible through the $\Z^2$-periodisation.  Fourth, the fixed-point
clause is stated translate-wise --
\texttt{literalFlow\_eq\_self\_of\_forall\_not\_mem\_bufferDisk} assumes that no integer
translate of the representative lies in a buffer disc -- because the torus flow is the
$\Z^2$-periodisation of the plane flow and the naive form would not suffice; points on the
boundaries are covered because \texttt{bufferDisk} is open, and the accumulation circle is
\texttt{literalFlow\_eq\_self\_of\_accumulationPoint}, no translate of
$p_\infty=(0,\tfrac12)$ lying in a buffer disc.  The two centre speeds are
\texttt{cellAngularSpeed\_cellCenter} ($\omega_n(0)=\Omega$) and
\texttt{cellAxialSpeed\_cellCenter} ($g_n(0)=\ell_nU+a_n$).

The uniform bound across the accumulating cells is carried out in
\texttt{Dynamics/Flow/PlaneFlowLipschitz.lean}, where a sum over cells is replaced by a
supremum: the buffer discs are disjoint and the flow is the identity off their union.  This
is not a deviation: the paper argues the same way in the proof of Proposition~4.1: the
natural lift of $\Phi_t$ is globally Lipschitz by
Lemma~\ref{lem:global-velocity}, its derivative is the cell formula on each disc and the
identity almost everywhere outside their union.

\begin{proof}
  \leanok
  \uses{def:dynamics:particle-flow, def:dynamics:topological-entropy, def:velocity:global-velocity, lem:global-velocity}
  On the $n $-th disc the velocity is
  $u(y,x_3)=(\omega_n(r)(y-p_n)^\perp,g_n(r)) $ with $r=\abs{y-p_n} $, so
  $\frac{\dd}{\dd t}\abs{y-p_n}^2=2\omega_n(r)(y-p_n)\cdot(y-p_n)^\perp=0 $: the radius is
  a constant of the motion, and the angular and axial coordinates move at the constant
  speeds $\omega_n(r) $ and $g_n(r) $.  This gives the closed formula, shows that no
  trajectory crosses a cell boundary, and matches the identity smoothly near
  $\partial D(p_n,R_n) $.  Differentiating the formula and using the two uniform shear
  bounds $\sup_n r\abs{\omega_n'(r)}\leq\abs\Omega\sup_s s\abs{\chi_h'(s)} $ and
  $\sup_n\abs{g_n'(r)}\leq\abs U\norm{\chi_h'}_\infty+\eta_{\rm ax} $ gives
  $\norm{D\Phi_t}\leq C(1+\abs t) $ on each cell with $C $ independent of $n $.  That the
  accumulation of the cells costs nothing is Lemma~\ref{lem:global-velocity}: the periodic
  lift of $\Phi_t $ is globally Lipschitz, its derivative is given by the cell formula on
  each disc and equals the identity almost everywhere outside their union, so the uniform
  cell bound is a bound on $\operatorname{Lip}(\Phi_t) $.  Applying the same estimate at
  $-t $ and using $\Phi_t^{-1}=\Phi_{-t} $ gives the two-sided bound after increasing the
  constant, and dividing its logarithm by $\abs t $ gives the vanishing log-Lipschitz rate.
  For the entropy, fix $s\in\R $ and an integer $m\geq1 $; the entropy formula for iterates
  and the Lipschitz entropy bound
  [Hasselblatt--Katok, Section~2.5f(3) and Proposition~2.5.5] give
  \[
    m\,h_{\mathrm{top}}(\Phi_s)
    =h_{\mathrm{top}}(\Phi_{ms})
    \leq3\log\bigl(C_{\rm L}(1+m\abs s)\bigr),
  \]
  and dividing by $m $ and letting $m\to\infty $ gives $h_{\mathrm{top}}(\Phi_s)=0 $.
\end{proof}

\section{The ideal induction group and the growth gap}

\begin{definition}[Ideal induction group]
  \label{def:dynamics:ideal-induction-group}
  \lean{FastDynamo.IsIdealInductionGroup, FastDynamo.HasVanishingIdealGrowthRate,
    FastDynamo.HasVanishingIdealOperatorGrowthRate,
    FastDynamo.ScaleFamily.TorusCellPlacement.idealTransport,
    FastDynamo.ScaleFamily.TorusCellPlacement.idealGroup,
    FastDynamo.mainIdealGroup}
  \leanok
  \uses{def:dynamics:particle-flow, def:velocity:global-velocity}
  The \emph{ideal induction equation} for the global velocity $u $ is
  \[
    \partial_tB+(u\cdot\nabla)B-(B\cdot\nabla)u=0,
    \qquad\Div B=0,
  \]
  that is, the induction equation of the introduction at $\eps=0 $.  Its solution group
  $\mathcal T_0(t) $ on $\Lp^2_\sigma(\T^3) $ is defined by the Cauchy formula
  \begin{equation}
    \bigl(\mathcal T_0(t)B_{\rm in}\bigr)(\Phi_t(a))
    =D\Phi_t(a)B_{\rm in}(a)
    \qquad\text{for a.e. }a,
    \label{eq:dynamics:ideal-cauchy}
  \end{equation}
  the magnetic field being transported and deformed by the particle flow.  Its
  exponential growth rate in operator norm is
  \begin{equation}
    \gamma_{\rm id}(u)
    :=\lim_{t\to\infty}\frac1t
    \log\norm{\mathcal T_0(t)}_{\Lp^2_\sigma\to\Lp^2_\sigma}.
    \label{eq:dynamics:ideal-rate}
  \end{equation}
\end{definition}

In the paper, $\mathcal T_0(t) $ is introduced in Section~1.2 as ``the solution group on
$\Lp^2_\sigma(\T^3) $ of the ideal induction equation''; the Cauchy formula is
\texttt{eq:ideal-cauchy-formula} of the \emph{proof} of Corollary~4.2, where it is derived from
the Lipschitz flow and shown to give that group, and $\gamma_{\rm id} $ is defined in
\texttt{eq:zero-ideal-growth} of the statement of Corollary~4.2.  The definition above takes
the Cauchy formula as the definition, which is the Lean model.

Stated \emph{and} constructed.  \texttt{IsIdealInductionGroup v} $\Phi$ $D\Phi$ $T$ bundles the
$\C_0$-group $T$ with the Cauchy representation \eqref{eq:dynamics:ideal-cauchy}, the
almost-everywhere Fr\'echet differentiability of a periodic lift of $\Phi_t$ (Rademacher),
the measure preservation of $\Phi_t$, and the two fields \texttt{solutionGroup} and
\texttt{unique}: $T$ is a solution group of the ideal equation of $v$ in the weak sense of
Definition~\ref{def:dynamics:ideal-weak-form}, and every solution group of that equation is
$T$.  Measure preservation replaces Liouville's formula, whose statement would need
$\det D\Phi_t$; with \texttt{unique} the group is determined by $v$ alone, which is the
paper's ``\emph{the} solution group'' of §1.2.

The group itself is built in two steps.  \texttt{idealTransport}
(\texttt{FastDynamo/Dynamics/Ideal/CauchyOperator.lean}) is
\eqref{eq:dynamics:ideal-cauchy} read as an operator on all of $\Lp^2(\T^3;\C^3)$, solved for
$\mathcal T_0(t)B$: the substitution $B\mapsto B\circ\Phi_{-t}$, which is a linear isometry of
$\Lp^2$ because $\Phi_{-t}$ preserves the volume, followed by pointwise multiplication by the
$\Lp^\infty$ matrix symbol $D\Phi_t(\Phi_{-t}\,\cdot)$ (the complexification
\texttt{flowFDerivC}, of norm at most $24(1+\abs t\,s)$ almost everywhere).  Then
\texttt{idealGroup} (\texttt{FastDynamo/Dynamics/Ideal/SolenoidalInvariance.lean})
corestricts it to $\Lp^2_\sigma$, which is the substantial step: see
Corollary~\ref{cor:ideal-growth}.  \texttt{mainIdealGroup} is the instance at the constructed
velocity.

The state space $\Lp^2_\sigma(\T^3) $ is \texttt{literalSolenoidalSubspace}
(\texttt{FastDynamo/Dynamics/Assembly/MainGrowthRate.lean}), which is \emph{defined} as the
image, under the isometric transport \texttt{toLiteralC3L2}, of the Fourier-defined solenoidal
subspace \texttt{productSolenoidalSubspace} of the spectral half.  Its identification with an
intrinsically defined subspace $\{B\in\Lp^2(\T^3;\C^3):\Div B=0\} $ of the literal torus
model is not formalised, no differential operator being transported to the literal torus.

The weak form of the ideal equation and the variational equation
\eqref{eq:dynamics:variational} (the paper's \texttt{eq:flow-variational-equation}), which the
paper uses in the proof of Corollary~4.2, are the nodes that follow.  Both need $\nabla u$ on
the literal torus, which is the derivative of the periodic lift of $u$ at every point
(Lemma~\ref{lem:dynamics:velocity-gradient}).  Remark~1.7 of the paper, the
vanishing-diffusivity limit, is Theorem~\ref{thm:vanishing-diffusivity} at the end of this
chapter.

\begin{definition}[The weak form of the ideal induction equation]
  \label{def:dynamics:ideal-weak-form}
  \lean{FastDynamo.literalFDeriv, FastDynamo.literalAdvection, FastDynamo.literalStretch,
    FastDynamo.literalFDeriv_literalTorusPoint, FastDynamo.literalAdvection_literalTorusPoint,
    FastDynamo.IsIdealTestField, FastDynamo.IsIdealWeakSolution,
    FastDynamo.IsIdealSolutionGroup, FastDynamo.IdealWeakIdentity}
  \leanok
  \uses{def:dynamics:ideal-induction-group}
  For a velocity $v$ on $\T^3$ whose periodic lift is differentiable, write $\nabla v$ for the
  derivative of the lift, $(v\cdot\nabla)\varphi$ and $(B\cdot\nabla)v=\nabla v\,B$.  A curve
  $W\in\C([-T,T];\Lp^2(\T^3;\CC^3))$ is a \emph{weak solution} of the ideal induction
  equation $\partial_tB+(v\cdot\nabla)B-(B\cdot\nabla)v=0$ on $[-T,T]$ if for every smooth
  periodic test field $\varphi$ and all $s,t\in[-T,T]$
  \[
    \inner{W(t)}{\varphi}-\inner{W(s)}{\varphi}
    =\int_s^t\Bigl(\inner{(W(\tau)\cdot\nabla)v}{\varphi}
      +\inner{W(\tau)}{(v\cdot\nabla)\varphi}\Bigr)\dd\tau .
  \]
  A \emph{solution group} is a one-parameter group $G(t)$ of bounded operators on
  $\Lp^2_\sigma(\T^3)$ all of whose orbits $t\mapsto G(t)B$ are weak solutions on every
  $[-T,T]$.
\end{definition}

The paper (proof of Corollary~4.2) understands solutions in the distributional sense in
$\C([-T,T];\Lp^2)$ and writes the weak form as
$\frac{\dd}{\dd t}\inner{B}{\varphi}=\inner{(B\cdot\nabla)u}{\varphi}+\inner{B}{(u\cdot\nabla)\varphi}$
for almost every $t$.  The Lean notion (\texttt{FastDynamo/Dynamics/IdealWeakForm.lean})
differs in these respects: the identity is tested against smooth spatial fields and integrated
in time (the equivalence with space--time test functions is not formalised); $\Div B=0$ is
not part of the notion of solution, only of the state space of a solution group; $\nabla v$
is the pointwise derivative of the lift; the integrals are Bochner integrals; and pairings are
$\int\langle\varphi,B\rangle$, conjugate-linear in the test field.  A solution group is
required to satisfy the group law, which the paper's ``solution group'' includes.

\begin{lemma}[The gradient of the constructed velocity]
  \label{lem:dynamics:velocity-gradient}
  \lean{FastDynamo.ScaleFamily.TorusCellPlacement.FullVelocityConcrete.fderiv_parameterizedFullGlobalPeriodizedRealSource_apply_pairBasis,
    FastDynamo.ScaleFamily.TorusCellPlacement.FullVelocityConcrete.transverseW1InfDeriv_parameterizedFullGlobalW1Inf_ae_eq_fderiv,
    FastDynamo.ScaleFamily.TorusCellPlacement.FullVelocityConcrete.fderiv_parameterizedFullGlobalPeriodizedRealSource_transverseDivergence,
    FastDynamo.ScaleFamily.TorusCellPlacement.hasFDerivAt_literalVelocity_lift,
    FastDynamo.ScaleFamily.TorusCellPlacement.norm_literalFDeriv_literalVelocity_le,
    FastDynamo.ScaleFamily.TorusCellPlacement.sum_literalFDeriv_literalVelocity_diag,
    FastDynamo.ScaleFamily.TorusCellPlacement.continuous_literalFDeriv_literalVelocity_literalFlow,
    FastDynamo.literalStretchL2,
    FastDynamo.ScaleFamily.TorusCellPlacement.literalStretchL2_coe_ae,
    FastDynamo.ScaleFamily.TorusCellPlacement.norm_literalStretchL2_le}
  \leanok
  \uses{def:velocity:global-velocity, lem:global-velocity, lem:main:velocity-regularity,
    def:dynamics:ideal-weak-form}
  The periodic lift of the global velocity $u$ is differentiable at every point of $\R^3$,
  with derivative bounded uniformly and divergence $\sum_i\partial_iu_i=0$ at every point; its
  transverse derivatives are, almost everywhere, the weak derivatives of $u$ as an element of
  $\W^{1,\infty}$; $t\mapsto\nabla u(\Phi_t(x))$ is continuous for every $x$; and
  $B\mapsto(B\cdot\nabla)u$ is a bounded operator on $\Lp^2(\T^3;\CC^3)$.
\end{lemma}

The paper's $\nabla u$ is the weak gradient of a $\W^{1,\infty}$ field and $\Div u=0$ holds
almost everywhere (Lemma~3.1 of the paper); here both hold at every point of the plane for the
closed-form lift, and the identification with the weak derivatives of the Fourier-graph field
\texttt{TransverseRealW1Inf} is the second item.  Modules
\texttt{FastDynamo/Global/TorusGlobalVelocityDerivativeBridge.lean} and
\texttt{FastDynamo/Dynamics/Ideal/VelocityGradient.lean}.

\begin{proof}
  \leanok
  \uses{lem:global-velocity, lem:main:velocity-regularity}
  Reduce a point to its representative in the unit box.  In a buffer disc the periodised field
  coincides near the point with the field of that cell, whose derivative is the derivative of
  the rescaled seed copy plus the axial offset; outside every buffer the derivative vanishes.
  The divergence vanishes cell by cell, for the helical part and for the axial offset.  The
  bound is the explicit derivative bound of the global velocity, and the almost-everywhere
  identification follows from the representative of the weak derivatives through the
  measure-preserving quotient map.
\end{proof}

\begin{lemma}[The variational equation]
  \label{lem:dynamics:variational-equation}
  \lean{FastDynamo.flowDeriv_apply_eq_add_integral,
    FastDynamo.ScaleFamily.TorusCellPlacement.flowFDeriv_eq_one_add_integral,
    FastDynamo.ScaleFamily.TorusCellPlacement.hasDerivAt_flowFDeriv,
    FastDynamo.ScaleFamily.TorusCellPlacement.flowFDeriv_apply_literalVelocity}
  \leanok
  \uses{def:dynamics:particle-flow, prop:particle-flow, lem:dynamics:velocity-gradient}
  For the particle flow $\Phi_t$ of the global velocity and its derivative $D\Phi_t$,
  \[
    D\Phi_t(a)=\id+\int_0^t(\nabla u)(\Phi_s(a))D\Phi_s(a)\dd s
  \]
  for every $a$ and every $t\in\R$; in particular $t\mapsto D\Phi_t(a)$ is differentiable with
  derivative $(\nabla u)(\Phi_t(a))D\Phi_t(a)$, and $D\Phi_t(a)u(a)=u(\Phi_t(a))$.
\end{lemma}

This is the paper's \texttt{eq:flow-variational-equation}, stronger: it holds at every point
and every time, where the paper has almost every $a$.  The route differs: the paper obtains it
from Cauchy--Lipschitz, Rademacher and the chain rule [Ambrosio--Crippa, Section~2]; here the
integral form of the flow equation is differentiated in the initial point under the integral
sign (dominated convergence and mean-value bounds), using the everywhere Fr\'echet derivative
of the closed-form lift.  The last identity, which the paper does not state, comes from
differentiating the group law of the lift.  Module
\texttt{FastDynamo/Dynamics/Flow/VariationalEquation.lean}.

\begin{proof}
  \leanok
  \uses{prop:particle-flow, lem:dynamics:velocity-gradient}
  Write $\Phi_t(a)=a+\int_0^tu(\Phi_s(a))\dd s$ on the lift.  The integrand is differentiable
  in $a$ with derivative $(\nabla u)(\Phi_s(a))D\Phi_s(a)$, bounded uniformly for $s$ in a
  compact interval, so differentiation under the integral sign gives the identity at every
  point.  The integrand is continuous in time
  (Lemma~\ref{lem:dynamics:velocity-gradient}), so the fundamental theorem of calculus gives
  the derivative.
\end{proof}

\begin{lemma}[The Cauchy formula solves the ideal equation]
  \label{lem:dynamics:ideal-weak-existence}
  \lean{FastDynamo.ScaleFamily.TorusCellPlacement.hasDerivAt_inner_literalFlow_flowFDerivC,
    FastDynamo.ScaleFamily.TorusCellPlacement.hasDerivAt_integral_inner_idealTransport,
    FastDynamo.ScaleFamily.TorusCellPlacement.isIdealWeakSolution_idealTransport,
    FastDynamo.ScaleFamily.TorusCellPlacement.isIdealSolutionGroup_idealGroup,
    FastDynamo.isIdealSolutionGroup_mainIdealGroup}
  \leanok
  \uses{def:dynamics:ideal-weak-form, def:dynamics:ideal-induction-group,
    lem:dynamics:variational-equation, lem:dynamics:velocity-gradient}
  For every $B\in\Lp^2(\T^3;\CC^3)$ and every $T$, the curve $t\mapsto\mathcal T_0(t)B$ given
  by the Cauchy formula \eqref{eq:dynamics:ideal-cauchy} is a weak solution of the ideal
  induction equation on $[-T,T]$, and
  \[
    \frac{\dd}{\dd t}\inner{\mathcal T_0(t)B}{\varphi}
    =\inner{(\mathcal T_0(t)B\cdot\nabla)u}{\varphi}
    +\inner{\mathcal T_0(t)B}{(u\cdot\nabla)\varphi}
  \]
  at every $t$.  Hence $\mathcal T_0$ is a solution group.
\end{lemma}

This is the paragraph ``Formula \texttt{eq:ideal-cauchy-formula} solves the ideal induction
equation'' of the paper's proof of Corollary~4.2, stronger: for every $B\in\Lp^2$ (not only
$\Lp^2_\sigma$), every $T$, and with the derivative at every time where the paper has almost
every $t$.  The route is the paper's: change of variables $x=\Phi_t(a)$, differentiation under
the integral, and the variational equation.  Module
\texttt{FastDynamo/Dynamics/Ideal/WeakSolution.lean}.

\begin{proof}
  \leanok
  \uses{lem:dynamics:variational-equation, lem:dynamics:velocity-gradient}
  By measure preservation and the Cauchy formula,
  $\inner{\mathcal T_0(t)B}{\varphi}=\int\langle\varphi(\Phi_t(a)),D\Phi_t(a)B(a)\rangle\dd a$.
  The integrand is differentiable in $t$ at every point, by the chain rule on the smooth lift
  of $\varphi$ and Lemma~\ref{lem:dynamics:variational-equation}, with a derivative bounded by
  $C\abs{B(a)}$ near each $t_0$; differentiation under the integral and the change of
  variables back give the displayed derivative, which is continuous in $t$.  The fundamental
  theorem of calculus gives the integrated identity.
\end{proof}

\begin{lemma}[Lipschitz test fields]
  \label{lem:dynamics:lipschitz-test-fields}
  \lean{FastDynamo.dense_idealTestFieldL2,
    FastDynamo.ScaleFamily.TorusCellPlacement.idealWeakIdentity_of_isIdealTestField,
    FastDynamo.IdealWeakIdentity.of_tendsto,
    FastDynamo.ScaleFamily.TorusCellPlacement.IsIdealWeakSolution.weakIdentity_of_lipschitz}
  \leanok
  \uses{def:dynamics:ideal-weak-form, lem:dynamics:velocity-gradient}
  The smooth periodic test fields are dense in $\Lp^2(\T^3;\CC^3)$.  If $W$ is a weak
  solution of the ideal equation of the global velocity on $[-T,T]$, the weak identity holds
  also for every test field $\varphi$ whose periodic lift is Lipschitz and differentiable at
  every point, with $\varphi,(u\cdot\nabla)\varphi\in\Lp^2$.
\end{lemma}

The paper extends identities from smooth to $\Hs^1$ test functions ``by density''.  Here the
extension is to Lipschitz, everywhere differentiable lifts, through Fourier partial sums:
integration by parts against the characters, Parseval convergence of the partial sums of
$\varphi$ and of $\partial_j\varphi$, and multiplication by the bounded components of $u$; no
mollifier is used.  The fields needed are $\psi\circ\Phi_{-\sigma}$ with $\psi$ smooth, which
are of this kind.  Module \texttt{FastDynamo/Dynamics/Ideal/LipschitzTestFields.lean}.

\begin{proof}
  \leanok
  \uses{lem:dynamics:velocity-gradient}
  Trigonometric polynomials are smooth and dense.  For a Lipschitz field with everywhere
  differentiable lift, the Fourier coefficients of $\partial_j\varphi$ are $2\pi im_j$ times
  those of $\varphi$ (integration by parts against characters), so the Fourier partial sums
  $\varphi_N$ converge to $\varphi$ and $(u\cdot\nabla)\varphi_N$ to $(u\cdot\nabla)\varphi$
  in $\Lp^2$.  The weak identity is linear and continuous in the test pair
  $(\varphi,(u\cdot\nabla)\varphi)$, so it passes to the limit.
\end{proof}

\begin{lemma}[Uniqueness of weak solutions; the solution group]
  \label{lem:dynamics:ideal-weak-uniqueness}
  \lean{FastDynamo.eq_zero_of_le_mul_abs_intervalIntegral,
    FastDynamo.ScaleFamily.TorusCellPlacement.inner_flowCompL2_eq_intervalIntegral,
    FastDynamo.ScaleFamily.TorusCellPlacement.norm_le_mul_abs_intervalIntegral,
    FastDynamo.ScaleFamily.TorusCellPlacement.IsIdealWeakSolution.eq_zero,
    FastDynamo.ScaleFamily.TorusCellPlacement.IsIdealSolutionGroup.eq_idealGroup,
    FastDynamo.ScaleFamily.TorusCellPlacement.existsUnique_isIdealSolutionGroup,
    FastDynamo.existsUnique_isIdealSolutionGroup_mainLiteralVelocity}
  \leanok
  \uses{def:dynamics:ideal-weak-form, lem:dynamics:ideal-weak-existence,
    lem:dynamics:lipschitz-test-fields, prop:particle-flow}
  If $W\in\C([-T,T];\Lp^2)$ is a weak solution of the ideal induction equation of the global
  velocity with $W(0)=0$, then $W=0$ on $[-T,T]$.  Consequently every solution group of that
  equation on $\Lp^2_\sigma(\T^3)$ is $\mathcal T_0$, and $\mathcal T_0$ is the unique solution
  group.
\end{lemma}

The first sentence is the paragraph ``Uniqueness follows from the same characteristic
argument'' of the paper's proof of Corollary~4.2.  The route differs: the paper obtains
$W(t,\Phi_t(a))=\int_0^t(\nabla u)(\Phi_s(a))W(s,\Phi_s(a))\dd s$ almost everywhere from
[Ambrosio--Crippa, Proposition~2.3] and applies Gronwall's inequality; here that identity is
proved in tested form,
$\inner{W(t)\circ\Phi_t}{\psi}=\int_0^t\inner{(W(\tau)\cdot\nabla)u}{\psi\circ\Phi_{-\tau}}\dd\tau$
for smooth $\psi$, directly from the weak form with the Lipschitz test fields
$\psi\circ\Phi_{-\sigma}$ (Lemma~\ref{lem:dynamics:lipschitz-test-fields}); density of the test
fields and the isometry $B\mapsto B\circ\Phi_t$ give
$\norm{W(t)}_2\leq\norm{S_u}\,\abs{\int_0^t\norm{W(\tau)}_2\dd\tau}$, with $S_u$ the bounded
operator $B\mapsto(B\cdot\nabla)u$ on $\Lp^2$, and Gronwall's inequality,
forward and backward in time, gives $W=0$.  Modules
\texttt{FastDynamo/Dynamics/Ideal/WeakSolutionUniqueness.lean} and
\texttt{FastDynamo/Dynamics/MainLagrangianClauses.lean}.

\begin{proof}
  \leanok
  \uses{lem:dynamics:lipschitz-test-fields, lem:dynamics:ideal-weak-existence}
  Test the weak identity of $W$ at $\psi\circ\Phi_{-\sigma}$ and differentiate the diagonal
  $\sigma\mapsto\inner{W(\sigma)}{\psi\circ\Phi_{-\sigma}}$; the transport terms cancel because
  $\frac{\dd}{\dd\sigma}\psi\circ\Phi_{-\sigma}=-(u\cdot\nabla)\psi\circ\Phi_{-\sigma}$, and what
  remains is the stretching term.  Integrating gives the tested characteristic identity; taking
  the supremum over unit $\psi$ gives the norm bound, and Gronwall's inequality gives $W=0$.
  For two solution groups $G$ and $\mathcal T_0$, the difference of their orbits from the same
  $B$ is a weak solution vanishing at $0$, so $G=\mathcal T_0$; existence is
  Lemma~\ref{lem:dynamics:ideal-weak-existence}.
\end{proof}

\begin{corollary}[No exponential growth for ideal induction]
  \label{cor:ideal-growth}
  \lean{FastDynamo.IsIdealInductionGroup, FastDynamo.MainLagrangianClauses,
    FastDynamo.HasVanishingIdealGrowthRate,
    FastDynamo.HasVanishingIdealOperatorGrowthRate,
    FastDynamo.MainLagrangianClauses.hasVanishingIdealGrowthRate,
    FastDynamo.MainLagrangianClauses.hasVanishingIdealOperatorGrowthRate,
    FastDynamo.inner_gradient_eq_zero_of_mem_literalSolenoidalSubspace,
    FastDynamo.ScaleFamily.TorusCellPlacement.isIdealInductionGroup_idealGroup,
    FastDynamo.ScaleFamily.TorusCellPlacement.norm_idealGroup_bounds,
    FastDynamo.ScaleFamily.TorusCellPlacement.tendsto_inv_mul_log_norm_idealGroup,
    FastDynamo.ScaleFamily.TorusCellPlacement.existsIdealGroupClauses,
    FastDynamo.exists_ne_zero_mainIdealGroup_rate,
    FastDynamo.existsUnique_isIdealSolutionGroup_mainLiteralVelocity}
  \leanok
  \uses{def:dynamics:ideal-induction-group, def:dynamics:particle-flow, def:velocity:global-velocity,
    def:dynamics:ideal-weak-form}
  Let $\mathcal T_0(t) $ be the solution group on $\Lp^2_\sigma(\T^3) $ of the ideal
  induction equation for the global velocity $u $ (a strongly continuous group of bounded
  operators, given by the Cauchy formula of
  Definition~\ref{def:dynamics:ideal-induction-group}).  Then, with $C_{\rm L} $ the
  constant of Proposition~\ref{prop:particle-flow},
  \[
    \frac1{C_{\rm L}(1+\abs t)}
    \leq\norm{\mathcal T_0(t)}_{\Lp^2_\sigma\to\Lp^2_\sigma}
    \leq C_{\rm L}(1+\abs t)
    \qquad(t\in\R).
  \]
  Consequently $\gamma_{\rm id}(u)=0 $.  Moreover, for every
  $B_{\rm in}\in\Lp^2_\sigma(\T^3) $ and $t\in\R $,
  \[
    \frac{\norm{B_{\rm in}}_2}{C_{\rm L}(1+\abs t)}
    \leq\norm{\mathcal T_0(t)B_{\rm in}}_2
    \leq C_{\rm L}(1+\abs t)\norm{B_{\rm in}}_2 ;
  \]
  in particular every non-zero ideal solution has exponential growth rate zero.
\end{corollary}

Every bound and rate of the statement is proved, for $\mathcal T_0$ the unique solution group
of the ideal equation in the weak sense of Definition~\ref{def:dynamics:ideal-weak-form}: it is
the group built from the Cauchy formula, it solves the equation
(Lemma~\ref{lem:dynamics:ideal-weak-existence}), and every solution group equals it
(Lemma~\ref{lem:dynamics:ideal-weak-uniqueness}); the fields \texttt{solutionGroup} and
\texttt{unique} of \texttt{IsIdealInductionGroup} record both, and
\texttt{existsUnique\_isIdealSolutionGroup\_mainLiteralVelocity} states them as one existence
and uniqueness statement.  The two-sided bounds are proved in the
\emph{per-datum} form, the display after \texttt{eq:zero-ideal-growth} in the statement,
$\norm B/(C_{\rm L}(1+\abs t))\leq\norm{T_tB}\leq C_{\rm L}(1+\abs t)\norm B$
-- \texttt{norm\_idealGroup\_bounds}, a conjunct of \texttt{MainLagrangianClauses} -- from
which the operator-norm form \texttt{eq:ideal-group-growth} follows by taking the supremum
over unit data, which is how the paper's proof obtains it too.  Both vanishing rates follow
-- the \emph{per-datum} one
$t^{-1}\log(\norm{T_tB}_2/\norm B_2)\to0$ for every $B\neq0$
(\texttt{tendsto\_inv\_mul\_log\_norm\_idealGroup}, packaged as
\texttt{HasVanishingIdealGrowthRate}, the last sentence of the statement)
and the \emph{operator-norm} one $t^{-1}\log\norm{\mathcal T_0(t)}\to0$
(\texttt{HasVanishingIdealOperatorGrowthRate}, which is both \texttt{eq:zero-ideal-growth}
here and the second dynamical clause of Theorem~1.2 of the paper (\texttt{thm:main}),
obtained from the bounds by
\texttt{MainLagrangianClauses.hasVanishingIdealOperatorGrowthRate}, which is the paper's
own argument: taking the supremum over unit initial data gives two-sided polynomial bounds
on the ideal operator norm, and the rate is squeezed between the logarithms of two linear
functions -- the squeeze itself is the reusable
\texttt{tendsto\_inv\_mul\_log\_opNorm\_of\_two\_sided\_bounds}).
\texttt{existsIdealGroupClauses} is the whole last conjunct of
\texttt{MainLagrangianClauses} in one statement, at the single constant
$C_{\rm L}=\texttt{dynamicalCL}=24(1+s)$.

None of these clauses is vacuous.  Every one of them -- the Cauchy formula, the group laws,
the strong continuity, the two-sided bounds and the per-datum rate -- is quantified over
$B\in\Lp^2_\sigma(\T^3) $, and would hold word for word if that space were $0 $; it is not,
by \texttt{exists\_ne\_zero\_literalSolenoidalSubspace}
(Theorem~\ref{thm:dynamics:solenoidal-semigroup}), whose witness is the spectral half's own
unstable eigenfunction.  \texttt{exists\_ne\_zero\_mainIdealGroup\_rate}
(\texttt{FastDynamo/Dynamics/MainLagrangianClauses.lean}) records the conclusion next to the
closing statements: there is a $B\neq0 $ in the state space at which
$t^{-1}\log(\norm{\mathcal T_0(t)B}/\norm B)\to0 $.

\textbf{Deviation: the solution notion.}  The paper's proof of Corollary~4.2 shows that the
Cauchy formula solves the ideal equation in weak form and that a distributional solution
$W\in\C([-T,T];\Lp^2)$ with $W(0)=0$ vanishes, so that $\mathcal T_0$ is the solution group.
Lean proves the same two facts for the weak notion of
Definition~\ref{def:dynamics:ideal-weak-form}, whose differences from the paper's
distributional notion (spatial smooth test fields and time-integrated identity, Bochner
integrals, pointwise $\nabla u$) are recorded there; the uniqueness argument differs from the
paper's (Lemma~\ref{lem:dynamics:ideal-weak-uniqueness}).

The substantial step of the construction is the one the paper's proof settles in one display
(testing against $\nabla\psi$, moving $\Phi_t$ onto the test function, and extending to
$\psi\circ\Phi_t\in\Hs^1$ by density): that $\Lp^2_\sigma$ is \emph{invariant}.  Read on the
literal torus it is the orthogonality
$\int_{\T^3}\overline{\nabla w}\cdot B=0$ for every solenoidal $B$ and every $C^1$ periodic
test function $w$ (\texttt{inner\_gradient\_eq\_zero\_of\_mem\_literalSolenoidalSubspace},
\texttt{FastDynamo/Dynamics/Ideal/GradientOrthogonality.lean}), proved by integration by parts
against the characters $e_m$ and Parseval.  Applying it at $w=e_m\circ\Phi_t$ needs a $C^1$
lift of $\Phi_t$, which the accumulation circle forbids, so the invariance is proved first for
the \emph{truncated} flows $\Phi^{(N)}$ -- finitely many cells, hence a $C^\infty$ lift --
and then transferred by $\Lp^2$-convergence
$\mathcal T_0^{(N)}(t)B\to\mathcal T_0(t)B$ and closedness of $\Lp^2_\sigma$
(\texttt{FastDynamo/Dynamics/Ideal/Truncated*.lean}).

\begin{proof}
  \leanok
  \uses{prop:particle-flow, def:dynamics:ideal-induction-group, def:dynamics:particle-flow,
    lem:dynamics:variational-equation, lem:dynamics:ideal-weak-existence,
    lem:dynamics:ideal-weak-uniqueness}
  Work with the periodic lift of $u$ and the corresponding lift of $\Phi_t$.  Cauchy--Lipschitz, Rademacher's theorem and the chain rule
  give the variational equation
  \begin{equation}
    D\Phi_t(a)=\id+\int_0^t(\nabla u)(\Phi_s(a))D\Phi_s(a)\dd s
    \qquad\text{for almost every }a
    \label{eq:dynamics:variational}
  \end{equation}
  [Ambrosio--Crippa, Section 2]; Liouville's formula and $\Div u=0 $ give
  $\det D\Phi_t=1 $, so the bi-Lipschitz map $\Phi_t $ preserves Lebesgue measure by the
  area formula.  Measure preservation turns the Cauchy formula into a bounded operator
  with $\norm{\mathcal T_0(t)B_{\rm in}}_2\leq\operatorname{Lip}(\Phi_t)\norm{B_{\rm in}}_2 $;
  testing against $\nabla\psi $ and moving $\Phi_t $ onto the test function shows the
  divergence constraint is preserved; differentiating
  $\inner{\mathcal T_0(t)B_{\rm in}}{\varphi} $ and using the variational equation gives
  the weak form of the ideal equation.  Uniqueness follows from the same characteristic
  representation plus Gronwall's inequality forward and backward in time, and the chain
  rule $D\Phi_{t+s}(a)=D\Phi_t(\Phi_s(a))D\Phi_s(a) $ gives the group law
  $\mathcal T_0(t+s)=\mathcal T_0(t)\mathcal T_0(s) $, $\mathcal T_0(-t)=\mathcal T_0(t)^{-1} $.
  Strong continuity follows from $\norm{\Phi_t-\id}_\infty\leq\norm u_\infty\abs t $ and
  $\norm{D\Phi_t-\id}_\infty\leq\e^{\norm{\nabla u}_\infty\abs t}-1 $ on a dense set of
  smooth fields, together with the local operator bound.  The Cauchy formula and
  Proposition~\ref{prop:particle-flow} give the per-datum upper bound
  $\norm{\mathcal T_0(t)B_{\rm in}}_2\leq C_{\rm L}(1+\abs t)\norm{B_{\rm in}}_2 $, and
  $\mathcal T_0(-t)\mathcal T_0(t)=\id $ gives the matching per-datum lower bound
  $\norm{B_{\rm in}}_2\leq C_{\rm L}(1+\abs t)\norm{\mathcal T_0(t)B_{\rm in}}_2 $.  Taking
  the supremum over unit initial data turns them into the displayed operator bounds, and
  taking logarithms and dividing by $t $ gives $\gamma_{\rm id}(u)=0 $ and the per-datum
  rate.
\end{proof}

\textbf{How the Lean proof differs from this one.}  Several steps are replaced by cheaper ones
that the closed form makes available.  Cauchy--Lipschitz is not used at all (the flow is
written down); Rademacher is not used (the lift is differentiable at every point, so the
almost-everywhere clause is discharged by \texttt{Filter.Eventually.of\_forall}), and the
variational equation is proved at every point by differentiating the integral form of the flow
equation (Lemma~\ref{lem:dynamics:variational-equation}); and Liouville's formula is replaced
by a pointwise Jacobian computation on the plane, $\det D\Phi_t=1$ at every point from the
closed formula, followed by the change of variables and a skew-product argument for the axial
shear.  The group law comes from the group law of the flow, and strong continuity is proved at
$t=0$ and moved along the group, by density of \emph{continuous} data and dominated convergence
for the symbol $D\Phi_t(a) $ at fixed $a $
(\texttt{FastDynamo/Dynamics/Ideal/StrongContinuity.lean}), instead of the paper's smooth data
and the bounds $\norm{\Phi_t-\id}_\infty\leq\norm u_\infty\abs t $,
$\norm{D\Phi_t-\id}_{\Lp^\infty}\leq\e^{L\abs t}-1 $.  The weak form is derived as in the paper
(Lemma~\ref{lem:dynamics:ideal-weak-existence}); uniqueness is proved from the weak form in
tested form rather than through [Ambrosio--Crippa, Proposition~2.3]
(Lemma~\ref{lem:dynamics:ideal-weak-uniqueness}).  The upper bound uses the bound on the
symbol $D\Phi_t $ rather than $\operatorname{Lip}(\Phi_t) $, and the lower bound is
$1\leq\norm{\mathcal T_0(-t)}\norm{\mathcal T_0(t)}$ exactly as here.  The last step is the
paper's: the per-datum bounds come first, and the operator norm follows by the supremum over
unit data.

\begin{proposition}[Particle flow and ideal induction]
  \label{prop:main-lagrangian}
  \lean{FastDynamo.MainLagrangianClauses, FastDynamo.mainLiteralVelocity,
    FastDynamo.MainTheoremLagrangian,
    FastDynamo.ScaleFamily.TorusCellPlacement.mainLagrangianClauses_literalVelocity,
    FastDynamo.mainTheoremLagrangian_holds}
  \leanok
  \uses{thm:main, def:velocity:global-velocity, def:dynamics:particle-flow, def:dynamics:topological-entropy, def:dynamics:ideal-induction-group}
  The velocity field $u $ of Theorem~\ref{thm:main} can be chosen so that there is
  $C_{\rm L}\geq1 $ with
  \[
    \operatorname{Lip}(\Phi_t)+\operatorname{Lip}(\Phi_t^{-1})
    \leq C_{\rm L}(1+\abs t),
    \qquad
    h_{\mathrm{top}}(\Phi_t)=0
    \qquad(t\in\R),
  \]
  and so that the ideal induction equation defines a strongly continuous solution group
  $\mathcal T_0(t) $ on $\Lp^2_\sigma(\T^3) $ with
  \[
    \lim_{t\to\infty}\frac1t
    \log\norm{\mathcal T_0(t)}_{\Lp^2_\sigma\to\Lp^2_\sigma}=0.
  \]
\end{proposition}

Proved.  The paper has no proposition of this shape: what it has is
Proposition~\ref{prop:particle-flow}, Corollary~\ref{cor:ideal-growth}, and the last
sentence of Theorem~1.2 of the paper, which quotes them.  The node exists because it is the
packaging of those statements \emph{at the velocity that
Theorem~\ref{thm:main} produces}, not at an arbitrary one, and because that packaging is
exactly the Lean predicate
\texttt{MainTheoremLagrangian}~$:=$~\texttt{MainLagrangianClauses mainLiteralVelocity}.

The paper makes the difference between this bundle and the theorem's own dynamical half
visible, and the Lean statement layer mirrors it.
\texttt{MainTheoremDynamicalClauses}~$v$ is exactly the two dynamical clauses of
Theorem~1.2 of the paper: an
ideal induction group for $v$ together with $h_{\mathrm{top}}(\Phi_t)=0$ at every $t\in\R$
and $t^{-1}\log\norm{\mathcal T_0(t)}\to0$, the entropy clause carrying the continuity of
every $\Phi_t$ with it because the paper defines $h_{\mathrm{top}}$ for a continuous self-map
of a compact metric space (Section~1.2 of the paper); \texttt{MainTheoremDynamical} is that
predicate
at \texttt{mainLiteralVelocity}, discharged by
\texttt{mainTheoremDynamical\_holds}.  It is a projection of this node --
\texttt{MainLagrangianClauses.toMainTheoremDynamicalClauses} -- and the bundle keeps what
the theorem drops: the inverse law $\Phi_{-t}\circ\Phi_t=\id$, the linear Lipschitz
bound of Proposition~\ref{prop:particle-flow}, and the two-sided per-datum bounds of
Corollary~\ref{cor:ideal-growth}.  \texttt{MainTheoremComplete} therefore carries both, as
the fields \texttt{dynamical} and \texttt{lagrangian}.

\texttt{mainLagrangianClauses\_literalVelocity}
(\texttt{FastDynamo/Dynamics/MainLagrangianClauses.lean}) proves the predicate at the general
parameters $(h,P,\Omega,U,\tau)$ with $\tau\geq0$, and \texttt{mainTheoremLagrangian\_holds}
instantiates it at the main parameter choice, where $\tau\geq0$ is
\texttt{mainParameterChoice.tau\_nonneg}.  The single constant shared by the Lipschitz
clause and the ideal bounds is
$C_{\rm L}=\texttt{dynamicalCL}=6\,\texttt{flowCL}=24(1+s)$ with
$s=\texttt{flowLipschitzSlope}$, the sum of the horizontal and axial shear bounds of the cell
profiles; the factor $6$ is what the passage from the sup metric of $\T^3$ to the Euclidean
metric of $\R^3$ and the complexification of $D\Phi_t$ cost.  The normalisation remark of the
paper (the unlabelled Remark~1.6) notes that the conclusion is invariant under $u\mapsto u/M $,
since the flow and the ideal group of the rescaled field are those of $u $ evaluated at time
$t/M $; that is Lemma~\ref{lem:main:normalisation}, where the bundle is proved for $u/M$ with
the constant $(1+M^{-1})C_{\rm L}$.

\begin{proof}
  \leanok
  \uses{prop:particle-flow, cor:ideal-growth, def:velocity:global-velocity,
    lem:global-velocity}
  The velocity field constructed for Theorem~\ref{thm:main} is the global velocity of
  Lemma~\ref{lem:global-velocity}, so Proposition~\ref{prop:particle-flow} gives the
  Lipschitz bound and the vanishing of the topological entropy of every time map, and
  Corollary~\ref{cor:ideal-growth} gives the solution group and its zero exponential growth
  rate.
\end{proof}

\begin{theorem}[Theorem~\ref{thm:main} in full]
  \label{thm:main-complete}
  \lean{FastDynamo.MainTheoremLagrangian, FastDynamo.MainTheoremDynamical,
    FastDynamo.MainTheoremComplete,
    FastDynamo.mainTheoremComplete_of_lagrangian,
    FastDynamo.mainTheoremComplete_holds}
  \leanok
  \uses{thm:main-strengthened, prop:main-lagrangian, thm:main}
  The conjunction of the strengthened spectral statement
  (Theorem~\ref{thm:main-strengthened}) with the dynamical clauses
  (Proposition~\ref{prop:main-lagrangian}) is Theorem~\ref{thm:main} in full, and more.
  In Lean it is the structure \texttt{MainTheoremComplete}, inhabited by
  \texttt{mainTheoremComplete\_holds}.
\end{theorem}

The spectral half is
\texttt{main\_theorem\_strengthened} and the dynamical half
\texttt{mainTheoremLagrangian\_holds}, and
\texttt{mainTheoremComplete\_of\_lagrangian} is the implication that combines them.  The
structure has three fields, not two: \texttt{strengthened}, then
\texttt{dynamical}~$:$~\texttt{MainTheoremDynamical}, the theorem's own two clauses,
and \texttt{lagrangian}~$:$~\texttt{MainTheoremLagrangian}, the larger bundle from which the
first is projected.  The implication form exists because the
module that defines \texttt{mainLiteralVelocity} cannot import the flow without a cycle;
\texttt{FastDynamo/Dynamics/MainLagrangianClauses.lean} supplies the closed term.
\texttt{Dynamics/ParticleFlow.lean} is not a leaf of the development: the whole of
\texttt{Dynamics/Flow/*} and \texttt{Dynamics/Ideal/*} depends on it.

\begin{proof}
  \leanok
  \uses{thm:main-strengthened, prop:main-lagrangian}
  Immediate: the two halves are proved separately and the structure is their conjunction.
\end{proof}


\section{Vanishing diffusivity}

Remark~1.7 of the paper (unlabelled, the last remark of §1.2) relates the two halves of
Theorem~\ref{thm:main}: the resistive evolution converges to the ideal one as
$\eps\downarrow0$.  The paper derives it from the Trotter--Kato approximation theorem
[Engel--Nagel, Theorem~III.4.8] and the Friedrichs commutator lemma [DiPerna--Lions,
Lemma~II.1].  The formalisation proves it by a direct energy estimate against a mollified ideal
orbit; the first three nodes are its ingredients, none of which the paper states.

\begin{lemma}[The mollifier on the torus]
  \label{lem:vanishing:mollifier}
  \lean{FastDynamo.literalTranslate, FastDynamo.continuous_literalTranslate_apply,
    FastDynamo.mollifierKernel, FastDynamo.integral_mollifierKernel,
    FastDynamo.literalMollify, FastDynamo.literalMollifyDeriv, FastDynamo.mollifierSymbol,
    FastDynamo.inner_flatTorus3FourierBasis_literalMollify,
    FastDynamo.inner_flatTorus3FourierBasis_literalMollifyDeriv,
    FastDynamo.norm_literalMollify_le, FastDynamo.tendsto_literalMollify,
    FastDynamo.exists_bound_sq_mul_mollifierSymbol}
  \leanok
  Let $\rho_\delta(z)=\delta^{-3}\rho(z/\delta)$ with $\rho\geq0$ a smooth bump on $\R^3$
  supported in the unit ball, of integral one, and let $J_\delta f=\int\rho_\delta(z)f(\cdot-z)\dd z$
  on $\Lp^2(\T^3;\CC^3)$, and $\partial_jJ_\delta$ the same with $\partial_j\rho_\delta$.  Then
  $\norm{J_\delta}\leq1$, $J_\delta f\to f$ as $\delta\downarrow0$, $J_\delta$ and
  $\partial_jJ_\delta$ are Fourier multipliers, and for each $\delta$ the symbol of $J_\delta$
  decays like $\abs m^{-2}$.
\end{lemma}

The operators are Bochner integrals of translations of $\Lp^2$
(\texttt{FastDynamo/Analysis/LiteralTorusMollifier.lean}).  The decay of the symbol is what makes
$J_\delta$ map $\Lp^2$ into $\Hs^2$.

\begin{proof}
  \leanok
  Translations are isometries and strongly continuous, so the integral has norm at most
  $\int\rho_\delta=1$ and converges to $f$ as the kernel concentrates.  The Fourier coefficient
  of a translate is the coefficient times a phase, which gives the multipliers; integrating by
  parts twice in $z$ gives the decay of the symbol.
\end{proof}

\begin{lemma}[The Friedrichs commutator lemma on the torus]
  \label{lem:vanishing:friedrichs}
  \lean{FastDynamo.literalComponentMul, FastDynamo.friedrichsCommutator,
    FastDynamo.friedrichsCommutator_apply, FastDynamo.exists_norm_friedrichsCommutator_le,
    FastDynamo.sum_integral_fderiv_mollifierKernel_mul_eq_zero,
    FastDynamo.tendsto_friedrichsCommutator}
  \leanok
  \uses{lem:vanishing:mollifier}
  Let $v$ be an $L$-Lipschitz velocity on $\T^3$ and
  $K_\delta f=\sum_j\bigl(\partial_jJ_\delta(v_jf)-v_j\partial_jJ_\delta f\bigr)$.  Then
  $\norm{K_\delta}_{\Lp^2\to\Lp^2}\leq C$ uniformly in $\delta\in(0,1]$.  If moreover the
  periodic lift of $v$ is differentiable at every point with $\Div v=0$ at every point, then
  $K_\delta f\to0$ in $\Lp^2$ as $\delta\downarrow0$, uniformly for $f$ in a compact subset of
  $\Lp^2$.
\end{lemma}

DiPerna and Lions state the lemma on $\R^d$, for Sobolev velocities and $\Lp^p$ fields, with
convergence in $\Lp^1_{\rm loc}$; here it is proved on the torus for Lipschitz $v$ with
everywhere differentiable lift, in $\Lp^2$, uniformly on compact sets, which is the case the
energy estimate needs; the constructed velocity satisfies the hypotheses
(Lemma~\ref{lem:dynamics:velocity-gradient}).  Module
\texttt{FastDynamo/Analysis/LiteralFriedrichsCommutator.lean}.

\begin{proof}
  \leanok
  \uses{lem:vanishing:mollifier}
  $K_\delta f=\sum_j\int\partial_j\rho_\delta(z)(\tau_zv_j-v_j\tau_z)f\dd z$ and
  $\norm{\tau_zv_j-v_j\tau_z}\leq L\abs z$, so
  $\norm{K_\delta}\leq3L\int\abs z\abs{\nabla\rho_\delta(z)}\dd z$, which does not depend on
  $\delta$.  For a continuous field $g$, the divergence term
  $\sum_j\int\partial_j\rho_\delta(z)(v_j(x-z)-v_j(x))\dd z=\int\rho_\delta(z)(\Div v)(x-z)\dd z$
  vanishes, which leaves
  $K_\delta g(x)=\sum_j\int\partial_j\rho_\delta(z)(v_j(x-z)-v_j(x))(g(x-z)-g(x))\dd z$, of size
  $O(L\sup_{\abs z\leq\delta}\norm{g(\cdot-z)-g}_\infty)$.  Continuous fields are dense, and
  the uniform bound extends the convergence to all of $\Lp^2$, uniformly on compact sets.
\end{proof}

\begin{lemma}[The resistive operator on mollified fields]
  \label{lem:vanishing:bridge}
  \lean{FastDynamo.productMollifyH2, FastDynamo.productH2Embedding_productMollifyH2,
    FastDynamo.toLiteralC3L2_productInductionAction_productMollifyH2,
    FastDynamo.ScaleFamily.TorusCellPlacement.hasDerivAt_literalMollify_idealGroup,
    FastDynamo.hasDerivAt_literalMollify_mainIdealGroup,
    FastDynamo.literalInductionSemigroup_apply_coe}
  \leanok
  \uses{lem:vanishing:mollifier, def:main:induction-operator,
    lem:dynamics:ideal-weak-existence, lem:dynamics:velocity-gradient,
    thm:dynamics:solenoidal-semigroup}
  For every $\delta$, $J_\delta$ maps $\Lp^2(\T^3;\CC^3)$ continuously into $\Hs^2$, and on
  $J_\delta f$ the auxiliary induction operator of the constructed velocity
  (Definition~\ref{def:main:induction-operator}) reads
  \[
    \mathcal A_\eps J_\delta f=\eps\Delta J_\delta f-\sum_ju_j\,\partial_jJ_\delta f
    +(J_\delta f\cdot\nabla)u .
  \]
  For every $B\in\Lp^2_\sigma$ the curve $t\mapsto J_\delta\mathcal T_0(t)B$ is differentiable
  at every $t$, with derivative
  $J_\delta\bigl((\mathcal T_0(t)B\cdot\nabla)u\bigr)-\sum_j\partial_jJ_\delta\bigl(u_j\,\mathcal T_0(t)B\bigr)$.
\end{lemma}

The first display is stated for the auxiliary action $\mathcal A_\eps$ on the product model and
transported to the literal torus; the derivative of the mollified ideal orbit comes from the
weak form tested against the characters.  Module
\texttt{FastDynamo/Dynamics/VanishingDiffusivity/ProductBridge.lean}.

\begin{proof}
  \leanok
  \uses{lem:vanishing:mollifier, lem:dynamics:ideal-weak-existence}
  The Fourier certificate of $J_\delta f$ in $\Hs^2$ is $(1+4\pi^2\abs m^2)\sigma_\delta(m)\hat
  f(m)$, square summable by the decay of the symbol; the derivative identities hold coefficient
  by coefficient, the transport term is computed on the continuous representative of $u$, and the
  stretching term on its pointwise derivative.  For the ideal orbit, the weak form at the
  characters $e_m$, placed in each vector coordinate, gives
  $J_\delta\mathcal T_0(s)B=J_\delta B+\int_0^sF$ with $F$ continuous, and the fundamental
  theorem of calculus gives the derivative.
\end{proof}

\begin{theorem}[Vanishing diffusivity; Remark~1.7 of the paper]
  \label{thm:vanishing-diffusivity}
  \lean{FastDynamo.resistiveEvolution,
    FastDynamo.norm_productInductionSemigroup_sub_literalMollify_le,
    FastDynamo.norm_mainIdealGroup_sub_resistiveEvolution_le,
    FastDynamo.tendstoUniformlyOn_resistiveEvolution_mainIdealGroup}
  \leanok
  \uses{thm:main, def:dynamics:ideal-induction-group, def:main:growth-rate}
  For each fixed $B_{\rm in}\in\Lp^2_\sigma$, $\e^{tL_{\eps,u}}B_{\rm in}$ converges strongly in
  $\Lp^2$ to $\mathcal T_0(t)B_{\rm in}$ as $\eps\downarrow0$, uniformly on bounded time
  intervals in $[0,\infty)$.
\end{theorem}

The statement is the paper's, for the velocity $u$ of Theorem~\ref{thm:main}: in Lean,
$u$ is \texttt{mainVelocity}, read on the literal torus as \texttt{mainLiteralVelocity},
$\mathcal T_0$ is \texttt{mainIdealGroup}, and the uniformity is on $[0,T]$ for every $T$.  The
resistive evolution \texttt{resistiveEvolution} $\eps\,t\,B$ is
\texttt{literalInductionSemigroup} at $\eps>0$, the semigroup generated by $L_{\eps,u}$ on
$\Lp^2_\sigma$ (Theorem~\ref{thm:dynamics:solenoidal-generator}); it is extended by $0$ to
$\eps\leq0$ and clamps negative times to $0$, conventions that the limit along
$\eps\downarrow0$ on $[0,T]$ does not see.  The paper's remark follows Theorem~1.2 and
concerns its velocity and its group $\mathcal T_0$; the Lean statement is for that velocity
only, not for a general divergence-free Lipschitz velocity.  The route differs: not
Trotter--Kato, but the energy estimate below.  Module
\texttt{FastDynamo/Dynamics/VanishingDiffusivity/Limit.lean}.

\begin{proof}
  \leanok
  \uses{lem:vanishing:friedrichs, lem:vanishing:bridge, lem:vanishing:mollifier,
    lem:dynamics:product-dissipative, thm:dynamics:product-semigroup}
  Fix a mollification scale $\delta$ and start the resistive semigroup at $J_\delta B$; let $w$
  be its $\Hs^2$ orbit and $e(t)=w(t)-J_\delta\mathcal T_0(t)B$, so $e(0)=0$.  By
  Lemma~\ref{lem:vanishing:bridge}, $e$ solves the resistive equation up to the residual
  $\eps\Delta J_\delta\mathcal T_0(t)B+K_\delta\mathcal T_0(t)B
  +(S_uJ_\delta-J_\delta S_u)\mathcal T_0(t)B$, with $S_uB=(B\cdot\nabla)u$,
  and the dissipativity $\re\inner{e}{\mathcal A_\eps e}\leq M_u\norm e^2$ with Gronwall's
  inequality bounds $\norm{e(t)}$ by the residual times $\e^{(2M_u+1)T}$.  On the compact orbit
  $\{\mathcal T_0(t)B:t\in[0,T]\}$ the commutator $K_\delta$ and the mollification errors are
  small uniformly for small $\delta$ (Lemmas~\ref{lem:vanishing:friedrichs}
  and~\ref{lem:vanishing:mollifier}); for fixed $\delta$ the diffusion term is $O(\eps)$; and
  the initial error $\norm{J_\delta B-B}$ is propagated by
  $\norm{\e^{tL_{\eps,u}}}\leq\e^{M_ut}$, with $M_u$ independent of $\eps$.  Choose $\delta$
  first and then $\eps$.
\end{proof}
